\section*{الفصل \ref{ch:g11:seismic-structure} --- الموجات الزلزالية تكشف باطن الأرض}

\begin{solution}{exo:g11:seismic-structure:1}
\omterm{def:g11:seismic-structure:body-waves}{الموجة P} تضغط الصخر وتمطه، فيتحرك على امتداد اتجاه
الانتقال؛ أما \omterm{def:g11:seismic-structure:body-waves}{الموجة S} فتقصه، فيتحرك الصخر عموديًا على ذلك
الاتجاه. وتعبر \omterm{def:g11:seismic-structure:body-waves}{الموجات P} الأجسام الصلبة والسوائل؛ أما \omterm{def:g11:seismic-structure:body-waves}{الموجات S} فلا تعبر إلا الأجسام الصلبة. وتُسجَّل
\omterm{def:g11:seismic-structure:body-waves}{الموجة P} أولًا لأنها الأسرع بين الاثنتين.
\end{solution}

\begin{solution}{exo:g11:seismic-structure:2}
$40 \times 111.2 = \qty{4448}{km}$، أي نحو \qty{4450}{km}. و$8900/111.2
\approx \ang{80}$.
\end{solution}

\begin{solution}{exo:g11:seismic-structure:3}
\qty{456}{s} تساوي 7~دقائق و36~ثانية؛ و\qty{822}{s} تساوي 13~دقيقة
و42~ثانية. والتأخر S--P هو $822 - 456 = \qty{366}{s}$، أي 6~دقائق
و6~ثوانٍ.
\end{solution}

\begin{solution}{exo:g11:seismic-structure:4}
$\sin i_c = 6.0/8.0 = 0.75$، إذن $i_c \approx \ang{48.6}$.
\end{solution}

\begin{solution}{exo:g11:seismic-structure:5}
سمك اللب الخارجي $5150 - 2891 = \qty{2259}{km}$. ونصف قطر اللب
$6371 - 2891 = \qty{3480}{km}$، و$3480/6371 \approx 0.55$: أي أكثر بقليل
من نصف نصف قطر الأرض.
\end{solution}

\begin{solution}{exo:g11:seismic-structure:6}
6~دقائق و6~ثوانٍ تساوي \qty{366}{s}: فالزلزال على بعد \ang{40}، أي
$40 \times 111.2 \approx \qty{4450}{km}$. والتأخر الواحد لا يعطي إلا
المسافة: فقد يقع المركز السطحي في أي نقطة من دائرة بهذا نصف القطر حول
المحطة. ويلزم محطتان أخريان لإيجاد النقطة الوحيدة الواقعة على البعد
الصحيح من الثلاث كلها.
\end{solution}

\begin{solution}{exo:g11:seismic-structure:7}
$i_c \approx \ang{48.6}$ و$\cos i_c \approx 0.661$. الموجة المباشرة: $200/6.0
\approx \qty{33.3}{s}$. \omterm{ex:g11:seismic-structure:sp}{الموجة المنكسرة الحرجة}: $200/8.0 + 2 \times 35 \times 0.661/6.0
\approx 25.0 + 7.7 = \qty{32.7}{s}$. \omterm{ex:g11:seismic-structure:sp}{فالموجة المنكسرة الحرجة} تصل أولًا، أبكر
بنحو \qty{0.6}{s}.
\end{solution}

\begin{solution}{exo:g11:seismic-structure:8}
$\tan i_c \approx 1.134$: فالمسافة الحرجة $2 \times 35 \times 1.134
\approx \qty{79}{km}$. \omterm{prop:g11:seismic-structure:crossover}{ومسافة العبور} $2 \times 35 \times
\sqrt{14.0/2.0} = 70\sqrt 7 \approx \qty{185}{km}$. ومحطة
\cref{exo:g11:seismic-structure:7} تقع عند \qty{200}{km}، فيما وراء
\omterm{prop:g11:seismic-structure:crossover}{مسافة العبور}: \omterm{ex:g11:seismic-structure:sp}{فالموجة المنكسرة الحرجة} تصل أولًا، كما وجدنا.
\end{solution}

\begin{solution}{exo:g11:seismic-structure:9}
$\sin i_c = 6.5/8.0 = 0.8125$، و$i_c \approx \ang{54.3}$ و$\cos i_c \approx
0.583$. \omterm{prop:g11:seismic-structure:head-wave}{زمن التقاطع}: $2 \times 6 \times 0.583/6.5 \approx \qty{1.1}{s}$.
\omterm{prop:g11:seismic-structure:crossover}{مسافة العبور}: $12\sqrt{14.5/1.5} \approx \qty{37}{km}$. تحت
قارة تقع \omterm{prop:g11:seismic-structure:crossover}{مسافة العبور} قرب \qty{165}{km}، فتلتقط محطات تفصل بينها بضع عشرات من
الكيلومترات المستقيمين كليهما؛ أما تحت المحيط فتصل \omterm{ex:g11:seismic-structure:sp}{الموجة المنكسرة الحرجة}
أولًا منذ نحو \qty{37}{km}، ولا يتجاوز طول مستقيم الموجة المباشرة
بضع عشرات من الكيلومترات: فيجب أن تفصل بين المحطات بضعة كيلومترات
لرسمه.
\end{solution}

\begin{solution}{exo:g11:seismic-structure:10}
$\sin r = (8.06/13.7)\sin\ang{30} \approx 0.294$، إذن $r \approx \ang{17.1}$:
ينحرف الشعاع نحو العمودي، أي نحو مركز
الأرض. ولذلك يعبر اللب على مسار طويل قريب من وسطه،
ويخرج بعيدًا في الجهة الأخرى: فأشعة P التي تدخل اللب
تخرج أبعد بكثير عن الزلزال من الأشعة التي تخطئه بقليل.
\end{solution}

\begin{solution}{exo:g11:seismic-structure:11}
$12\,742/13.7 \approx \qty{930}{s}$، أي نحو 15.5~دقيقة، في حين تستغرق الموجة
\qty{1210}{s}، أي نحو 20~دقيقة: فهي متأخرة بنحو 280~ثانية. فحتى بأعلى
سرعة في الوشاح كانت الموجة ستصل أبكر بكثير، فلا بد أن
وسط الأرض يبطئها. وقد استنتج أولدهام أن للأرض لبًّا
مركزيًا تنتقل فيه \omterm{def:g11:seismic-structure:body-waves}{الموجات P} أبطأ منها في الوشاح.
\end{solution}

\begin{solution}{exo:g11:seismic-structure:12}
(1) بحسب المبرهنة، $t_i = 2h\cos i_c/v_1$، و$\cos i_c = \sqrt{1 -
v_1^2/v_2^2} = \sqrt{v_2^2 - v_1^2}/v_2$. إذن
\[
  t_i = \frac{2h\sqrt{v_2^2 - v_1^2}}{v_1 v_2},
  \qquad\text{ومن ثم}\qquad
  h = \frac{t_i\,v_1 v_2}{2\sqrt{v_2^2 - v_1^2}}.
\]
(2) $h = 6.23 \times 6.2 \times 8.1/\bigl(2\sqrt{65.61 -
38.44}\bigr) \approx 312.9/10.43 \approx \qty{30.0}{km}$.
\end{solution}

\begin{solution}{exo:g11:seismic-structure:13}
(1) $(3480/6371)^3 \approx 0.163$: فاللب نحو \qty{16}{\percent} من
الحجم. (2) متوسط الكثافة هو المتوسط المرجّح بالحجوم: $5.52 = 0.163
\rho_c + 0.837 \times 4.45$، إذن $\rho_c \approx (5.52 - 3.72)/0.163 \approx
\qty{11.0}{g/cm^3}$. (3) وهذا أكثر من \qty{9.9}{g/cm^3}، كما ينبغي
أن يكون، لأن كثافة اللب تزداد نحو المركز (حتى
\qty{13.1}{g/cm^3})؛ وهي نحو أربعة أضعاف كثافة القشرة،
\qty{2.6}{g/cm^3}. فاللب مصنوع من مادة أكثف بكثير من صخور
السطح.
\end{solution}

\begin{solution}{exo:g11:seismic-structure:14}
(1) $105 \times 111.2 \approx \qty{11700}{km}$ و$143 \times 111.2 \approx
\qty{15900}{km}$. (2) $(\cos\ang{105} - \cos\ang{143})/2 \approx (-0.259 +
0.799)/2 \approx 0.27$: أي نحو \qty{27}{\percent} من السطح، على هيئة حلقة حول
النقطة المقابلة للزلزال. (3) \omterm{def:g11:seismic-structure:shadow-zone}{منطقة الظل} مدى من
المسافات عن المركز السطحي: فهي تتحرك مع الزلزال. ولكل زلزال
حلقته الخاصة.
\end{solution}

\begin{solution}{exo:g11:seismic-structure:15}
(1) $h = 9.5 \times 6.3 \times 8.1/\bigl(2\sqrt{65.61 - 39.69}\bigr) \approx
484.8/10.18 \approx \qty{48}{km}$. (2) $\sin i_c = 6.3/8.1 \approx 0.778$،
و$\cos i_c \approx 0.629$، و$t_i = 2 \times 30 \times 0.629/6.3 \approx
\qty{6.0}{s}$. (3) سمك القشرة تحت السلسلة نحو \qty{48}{km}،
أي أسمك من المتوسط القاري البالغ نحو \qty{41}{km}، وقريب من
متوسط الأحزمة الجبلية: فالسلاسل العالية تقوم على قشرة سميكة.
\end{solution}

\begin{solution}{pb:g11:seismic-structure:1}
\textbf{1.} تقع نقاط Pg على مستقيم؛ وتقع نقاط Pn على
مستقيم آخر أقل انحدارًا.

\textbf{2.} الزلزال على عمق ضحل جدًا تحت المركز السطحي:
فالمحطة الواقعة على مسافة صفر كانت ستسجل الموجة فورًا.

\textbf{3.} $(32.3 - 3.3)/(200 - 20) = 29.0/180 \approx \num{0.161}$~\unit{s}
لكل كيلومتر.

\textbf{4.} $v_1 = 180/29.0 \approx \qty{6.2}{km/s}$.

\textbf{5.} $3.3 + 280 \times 0.161 \approx \qty{48.4}{s}$.

\textbf{6.} $(43.2 - 18.6)/(300 - 100) = 24.6/200 = \num{0.123}$~\unit{s}
لكل كيلومتر، إذن $v_2 = 200/24.6 \approx \qty{8.13}{km/s}$.

\textbf{7.} إنها قريبة من \qty{8.0}{km/s}: فتحت القشرة يقع صخر
أعلى الوشاح.

\textbf{8.} $\sin i_c = 6.21/8.13 \approx 0.763$، إذن $i_c \approx
\ang{49.8}$.

\textbf{9.} تهبط \omterm{ex:g11:seismic-structure:sp}{الموجة المنكسرة الحرجة} بالزاوية $i_c$ وتصعد
بالزاوية نفسها. وأول نقطة يمكنها أن تعود فيها إلى السطح تقع
على بعد $h\tan i_c$ وراء النقطة التي بلغت فيها \omterm{def:g11:seismic-structure:moho}{الموهو}، وهذه نفسها على بعد $h\tan i_c$
من المصدر: فلا تستطيع أي محطة أقرب من $2h\tan i_c$ أن تتلقاها.

\textbf{10.} يسجل مقياس الزلازل كل موجة تبلغه، لا الأولى
وحدها؛ فالموجة Pn تصل بعد الموجة Pg لكنها تظل مرئية موجةً منفصلة
على السجل.

\textbf{11.} عند \qty{160}{km} تسبق Pg (\qty{25.8}{s}) الموجة Pn
(\qty{26.0}{s})؛ وعند \qty{180}{km} تسبق Pn (\qty{28.5}{s}) الموجة Pg
(\qty{29.0}{s}). فالتحول يحدث بين هاتين المحطتين.

\textbf{12.} Pg: $t = 3.3 + 0.1611\,(x - 20) \approx 0.08 + 0.1611\,x$. وPn:
$t = 18.6 + 0.123\,(x - 100) = 6.30 + 0.123\,x$ ($x$ بالكيلومترات، و$t$
بالثواني).

\textbf{13.} $0.08 + 0.1611\,x = 6.30 + 0.123\,x$ يعطي $x \approx
6.22/0.0381 \approx \qty{163}{km}$.

\textbf{14.} الموجة Pn، بعد $6.30 + 0.123 \times 250 \approx \qty{37.1}{s}$.

\textbf{15.} $t_i \approx \qty{6.3}{s}$.

\textbf{16.} $\cos i_c \approx 0.646$، إذن $h = 6.21 \times 6.3/(2 \times
0.646) \approx \qty{30.3}{km}$.

\textbf{17.} $x_c = 2 \times 30.3 \times \sqrt{14.34/1.92} \approx 60.6 \times
2.73 \approx \qty{165}{km}$، أي في حدود \qty{2}{km} من قيمة \qty{163}{km} في
السؤال~13.

\textbf{18.} $\tan i_c \approx 1.18$: $2 \times 30.3 \times 1.18 \approx
\qty{72}{km}$. وأقرب محطة تسجل Pn تقع عند \qty{80}{km}، والمحطة
الواقعة عند \qty{60}{km} لا تسجلها: فالنتيجتان متفقتان.

\textbf{19.} تُقرأ الأزمنة بدقة \qty{0.1}{s}، والخطأ بمقدار
\qty{0.1}{s} في $t_i$ وحده يغيّر $h$ بنحو \qty{0.5}{km}؛ والسرعتان،
المأخوذتان من الميلين، غير يقينيتين أيضًا، ولم يكن الزلزال عند
السطح تمامًا. ثم إن القشرة ليست متجانسة تمامًا. فدقة كيلومتر
أو اثنين هي أقصى ما تسمح به البيانات.

\textbf{20.} \textbf{يقع \omterm{def:g11:seismic-structure:moho}{الموهو} على عمق نحو \qty{30}{km} تحت هذه المنطقة}:
فالقشرة هنا رقيقة نسبيًا بالنسبة إلى قارة، يبلغ متوسط سمك قشرتها نحو
\qty{41}{km}؛ و\qty{30}{km} سمك نموذجي للمناطق التي
مُطّت فيها القشرة.
\end{solution}
